% Vietnamese translation for OI Public Library
% Source-File: IOI中国国家候选队论文/2005/张伟达/张伟达.ppt
\documentclass[11pt]{article}
\usepackage{oipl-vn}

\newcommand{\Oh}{\mathcal{O}}

\title{Bản trình chiếu: Cải tiến thuật toán khi tăng quy mô và số chiều}
\author{Zhang Weida}
\date{2005}

\begin{document}
\maketitle

\origtitle{Dùng tư tưởng cải tiến thuật toán để giải bài tăng quy mô và số chiều: bản trình chiếu}
\sourcepath{IOI China candidate team papers / 2005 / Zhang Weida / Zhang Weida.ppt}

\section{Mạch chính của slide}

Bản trình chiếu đi kèm bài viết về tư tưởng cải tiến thuật toán khi bài toán
tăng quy mô hoặc tăng số chiều. Phần văn bản trích xuất được từ PowerPoint cũ
chủ yếu tập trung vào \textit{Team Selection} của Balkan OI 2004 Day 1, với
giới hạn \(3\le N\le 500000\), quan hệ ``\(A\) tốt hơn \(B\)'', khái niệm
\textit{excellent}, các vòng lặp \(X,Y\), đại lượng \(B[A_i]\), giá trị
\(\operatorname{best}_i\), cặp \((\mathit{key},\mathit{value})\), và truy vấn
giá trị nhỏ nhất trong đoạn khóa. Bản dịch này giữ luồng của slide và dùng
bài viết đã dịch làm nền để nối các bước bị mất chữ.

Thông điệp trung tâm là: khi một bài quen thuộc được mở rộng thêm chiều, không
nên chỉ thêm chỉ số vào công thức cũ. Cần hỏi thuật toán gốc thật sự dùng tính
chất nào, rồi đưa phần tăng thêm vào đúng chỗ bằng thứ tự duyệt, liệt kê,
tham lam hoặc cấu trúc dữ liệu.

\section{Từ mô phỏng trực tiếp tới tập ứng viên}

Bài \textit{Team Selection} có \(N\) thí sinh, mỗi thí sinh có thứ hạng trong
ba cuộc thi. Thí sinh \(A\) tốt hơn \(B\) nếu \(A\) xếp trước \(B\) ở cả ba
cuộc thi. Một thí sinh là \textit{excellent} nếu không có ai tốt hơn mình.
Yêu cầu là đếm số thí sinh excellent.

Thuật toán đầu tiên mà slide ghi lại là hai vòng lặp:
\[
  \text{for }X\quad \text{for }Y\quad \text{kiểm tra }Y\text{ tốt hơn }X.
\]
Mỗi phép kiểm tra chỉ cần so sánh ba thứ hạng, nhưng tổng thể là
\(\Oh(N^2)\), không thể chạy với \(N=500000\).

Bước cải tiến đầu tiên là duyệt theo thứ hạng cuộc thi thứ nhất. Khi xét
thí sinh \(X\), chỉ các thí sinh đứng trước \(X\) ở cuộc thi thứ nhất mới có
khả năng tốt hơn \(X\). Điều này làm rõ một chiều của bài toán nhưng chưa đủ
giảm độ phức tạp.

Bước tiếp theo dùng tập excellent đã biết. Nếu \(Y\) không excellent vì có
\(Z\) tốt hơn \(Y\), thì bất kỳ khi nào \(Y\) tốt hơn \(X\), \(Z\) cũng tốt
hơn \(X\). Vì vậy ta không cần giữ \(Y\) như một ứng viên riêng. Khi xét
\(X\), chỉ cần so sánh với các thí sinh excellent đã gặp. Độ phức tạp trở
thành \(\Oh(NK)\), với \(K\) là số thí sinh excellent, nhưng \(K\) vẫn có thể
lớn. Slide dùng bước này để minh họa rằng cải tiến cục bộ có ích cho tư duy,
nhưng chưa chắc đủ cho dữ liệu lớn.

\section{Hạ chiều: bài hai cuộc thi}

Để thấy bản chất, slide quay về bài chỉ có hai cuộc thi. Gọi \(A_i\) là thí
sinh đứng thứ \(i\) ở cuộc thi thứ nhất, và \(B[A_i]\) là thứ hạng của thí
sinh ấy ở cuộc thi thứ hai. Duyệt dãy \(A_1,A_2,\ldots,A_N\). Khi đang xét
\(A_i\), điều kiện ``xếp trước ở cuộc thi thứ nhất'' đã tự động đúng với mọi
\(A_j\), \(j<i\). Do đó \(A_i\) không excellent khi và chỉ khi
\[
  \exists j<i:\ B[A_j]<B[A_i].
\]

Ta chỉ cần duy trì giá trị tốt nhất đã gặp:
\[
  \operatorname{best}_i=\min_{j<i}B[A_j].
\]
Nếu \(\operatorname{best}_i<B[A_i]\), có người tốt hơn \(A_i\); ngược lại
\(A_i\) là excellent. Bài hai chiều được giải trong \(\Oh(N)\).

Đây là bước quan trọng của slide. Cải tiến không phải là thêm mẹo vào vòng
lặp, mà là chọn một thứ tự duyệt để một điều kiện tự biến mất, rồi nén điều
kiện còn lại thành một đại lượng duy nhất.

\section{Mở rộng lên ba cuộc thi}

Quay lại ba cuộc thi, đặt \(C[A_i]\) là thứ hạng của \(A_i\) ở cuộc thi thứ
ba. Khi xét \(X=A_i\), ta cần biết trong các thí sinh đã xử lý có ai thỏa
\[
  B[A_j]<B[A_i]\quad\text{và}\quad C[A_j]<C[A_i]
\]
hay không. Điều kiện thứ nhất vẫn được giải quyết bởi thứ tự duyệt. Hai điều
kiện còn lại tạo thành một bài truy vấn hai chiều: trong các khóa
\(B[A_j]\) nhỏ hơn \(B[A_i]\), tìm giá trị nhỏ nhất của \(C[A_j]\).

Slide biểu diễn mỗi thí sinh đã gặp bằng một cặp
\[
  (\mathit{key},\mathit{value})=(B[X],C[X]).
\]
Cấu trúc dữ liệu cần hỗ trợ:
\begin{itemize}
  \item chèn một cặp \((\mathit{key},\mathit{value})\);
  \item hỏi giá trị \(\mathit{value}\) nhỏ nhất trong các key nhỏ hơn key hiện
  tại.
\end{itemize}
Một cây đoạn trên miền thứ hạng của cuộc thi thứ hai làm đúng việc này. Mỗi
lá ứng với một key; mỗi nút lưu giá trị \(C\) nhỏ nhất trong đoạn key của nó.
Khi chèn, cập nhật các nút trên đường từ lá tới gốc. Khi hỏi, phân rã đoạn
\([1,\mathit{key}-1]\) thành các đoạn chuẩn và lấy min.

Thuật toán cuối cùng:
\begin{enumerate}
  \item duyệt thí sinh theo thứ tự cuộc thi thứ nhất;
  \item với thí sinh \(X\), hỏi \(\min C[Y]\) trong các \(Y\) đã gặp và có
  \(B[Y]<B[X]\);
  \item nếu giá trị này nhỏ hơn \(C[X]\), \(X\) bị một thí sinh khác trội hơn;
  \item ngược lại, \(X\) là excellent;
  \item chèn \((B[X],C[X])\) vào cây.
\end{enumerate}
Mỗi bước mất \(\Oh(\log N)\), nên tổng thời gian là \(\Oh(N\log N)\) và bộ
nhớ tuyến tính theo \(N\).

\section{Ý nghĩa của tư tưởng cải tiến}

Ví dụ \textit{Team Selection} thể hiện đủ các tầng cải tiến mà bài viết muốn
nhấn mạnh. Cách làm thô so sánh mọi cặp. Cách duyệt theo cuộc thi thứ nhất
loại bỏ một phần ứng viên. Cách giữ tập excellent giảm lặp lại nhưng vẫn chưa
đủ. Bài hai cuộc thi cho thấy điều kiện đã được bảo đảm bởi thứ tự duyệt có
thể bỏ khỏi trạng thái. Cuối cùng, bài ba cuộc thi giữ đúng ý tưởng ấy và thay
đại lượng ``tốt nhất đã gặp'' bằng truy vấn min trên một miền khóa.

\section{Kết luận}

Bản trình chiếu không nhằm giới thiệu riêng cây đoạn, mà dùng cây đoạn như
một kết quả tự nhiên của phân tích. Với bài toán tăng số chiều, việc thêm cấu
trúc dữ liệu chỉ nên đến sau khi ta đã biết chiều nào được xử lý bằng thứ tự
duyệt, chiều nào làm khóa, và chiều nào là giá trị cần tối ưu. Đó là điểm
khác nhau giữa ``cố áp dụng thuật toán cũ'' và thật sự cải tiến thuật toán.

\end{document}
